\documentclass{article}

\usepackage[a4paper,margin=2cm]{geometry}
\usepackage{../../../components/components}
\usepackage{colortbl}
\usepackage{array}
\usepackage{hyperref}
\usepackage{stmaryrd}

\usepackage{fancyhdr}
\pagestyle{fancy}
\fancyhf{}
\fancyhead[L]{DL2 Math-Info}
\fancyhead[C]{Révisions de Pâques}
\fancyhead[R]{2025-2026}
\fancyfoot[L]{Ewen Rodrigues de Oliveira}
\fancyfoot[R]{\thepage}

\begin{document}

\thispagestyle{empty}

\newcommand{\doctitle}{Révisions pour les examens}

\vspace*{\fill}

\begin{center}
    {\LARGE\textbf{\doctitle}}\\[0.6em]
    {\Large Double licence 2 Mathématiques et Informatique  Fondamentale}\\[2em]
    {\large Année universitaire 2025-2026}\\[0.25em]
    {\large par \textit{Ewen Rodrigues de Oliveira}}
\end{center}

\vfill

\noindent
\textbf{Le livret est une suite du livret de révision Pâques. Celui-ci est plus long et cible directement des annales d'examens, tandis que le livret de Pâques était plus court et ciblait des exercices de TD pour progresser sur les notions de base.}
\medskip

\noindent
Ce livret propose des révisions sur les matières d'Éléments d'algorithmique, de Probabilités Discrètes, ainsi que d'Analyse et d'Algèbre \textit{(et un peu de langage C)}. Je n'ai pas créé ces exercices : j'ai surtout rassemblé des exercices intéressants issus d'annales et de feuilles de TD de l'\href{https://u-pariscite.fr/}{Université Paris Cité}.

\medskip

\noindent
S'il y a une erreur, une coquille ou un énoncé ambigu, merci de me contacter pour que je puisse corriger cela rapidement. Un corrigé (non exhaustif) sera publié sur \href{https://ewenrdo.fr/ressources/}{mon site}.

\begin{flushright}
\textbf{Bonnes révisions \faSmile}
\end{flushright}

\newpage

\section*{Crédits}
\subsection*{Mathématiques}
\begin{itemize}
\item TD et annales de l'Université Paris Cité de 2020-2021 à 2024-2025.
\item Bibmath, section Math-Spé
\item Ministère de l'Éducation nationale, de la Jeunesse et des Sports, annales du CAPES de mathématiques 2025.
\item Sujet de contrôle continu de Pierre Faugere pour l'année 2025-2026.
\end{itemize}

\subsection*{Informatique}
\begin{itemize}
\item Cours et épreuves de Dominique Poulalhon à l'Université Paris Cité pour les années 2023-2024 et 2024-2025.
\item Sujet de projet de Juliusz Chroboczek en Langage C pour l'année 2025-2026.
\end{itemize}

\tableofcontents

\newpage
\renewcommand{\arraystretch}{1.3}
\section{Jour 1}
\subsection{Algèbre}
\exercise{1}{Étude d'un endomorphisme}{Soit $E$ un espace vectoriel euclidien orienté muni d'une base orthonormée directe $\mathcal{B} = (i, j, k)$. Soit $f \in \mathcal{L}(E)$ dont la matrice dans la base $\mathcal{B}$ est
\[ \frac{1}{2} \begin{pmatrix} 
1 & -\sqrt{2} & 1 \\
\sqrt{2} & 0 & -\sqrt{2} \\
1 & \sqrt{2} & 1 
\end{pmatrix} \]
\begin{enumerate}
    \item Former une base orthonormée directe $\mathcal{B}' = (u, v, w)$ telle que $v, w \in P = \{x + z = 0\}$.
    \item Former la matrice de $f$ dans $\mathcal{B}'$ et reconnaître $f$.
\end{enumerate}}

\exercise{2}{Inégalité de Rayleigh}{Soit $u$ un endomorphisme autoadjoint d'un espace euclidien $E$ de dimension $n$. On note $\lambda_1 \le \lambda_2 \le \dots \le \lambda_n$ les valeurs propres de $u$, comptées avec leur multiplicité. Démontrer que, pour tout $x \in E$, on a
\[ \lambda_1 \|x\|^2 \le \langle u(x), x \rangle \le \lambda_n \|x\|^2. \]}

\exercise{3}{Similitudes}{Soit $E$ un espace euclidien et $f \in \mathcal{L}(E)$ tel que
\[ \forall x, y \in E, x \perp y \Rightarrow f(x) \perp f(y) \]
Montrer qu'il existe $\lambda \in \mathbb{R}^+$ tel que
\[ \forall x \in E, \|f(x)\| = \lambda \|x\| \]
\textit{Indication : considérer une base orthonormée $(e_1, e_2, \dots, e_n)$, poser $\lambda_i = \|f(e_i)\|$, et calculer $\langle f(e_i + e_j), f(e_i - e_j) \rangle$}}

\subsection{Probabilités}
\exercise{1}{Chaîne à deux états}{On considère une chaîne de Markov à valeurs dans $E = \{0, 1\}$ de matrice de transition
\[ P = \begin{pmatrix} 1-p & p \\ q & 1-q \end{pmatrix}, \quad p, q \in [0, 1]. \]
On note $\mu_n = (P(X_n = 0), P(X_n = 1))$ la loi de $X_n$.
\begin{enumerate}
    \item[(a)] Montrer que si $\mu_0 = (\alpha, 1 - \alpha)$, alors
    \[ \mu_1 = (\alpha(1 - p) + (1 - \alpha)q, \alpha p + (1 - \alpha)(1 - q)). \]
    \item[(b)] Déterminer les distributions stationnaires.
    \item[(c)] On suppose $p > 0$ et $q > 0$. Montrer qu'il existe une unique distribution stationnaire.
    \item[(d)] Dans le cas particulier $p = q = 1/2$, calculer $\mu_1$ puis $\mu_2$ pour une loi initiale quelconque $\mu_0 = (\alpha, 1 - \alpha)$.
\end{enumerate}}

\exercise{2}{Étude d'une chaîne sur $\{0, 1, 2\}$}{On considère la chaîne de Markov sur $E = \{0, 1, 2\}$ de matrice
\[ P = \begin{pmatrix} 1/2 & 1/2 & 0 \\ 1/4 & 1/2 & 1/4 \\ 0 & 1/2 & 1/2 \end{pmatrix}. \]
\begin{enumerate}
    \item[(a)] Montrer que la chaîne est irréductible.
    \item[(b)] Montrer qu'elle est apériodique.
    \item[(c)] Déterminer une distribution stationnaire.
    \item[(d)] Calculer $\mu_1$ et $\mu_2$ lorsque $\mu_0 = (1, 0, 0)$.
\end{enumerate}}

\exercise{3}{Marche aléatoire sur un triangle}{On considère la marche aléatoire simple sur le triangle $\{A, B, C\}$ : à chaque instant, la particule saute uniformément vers l'un des deux autres sommets.
\begin{enumerate}
    \item[(a)] Écrire la matrice de transition.
    \item[(b)] La chaîne est-elle irréductible ?
    \item[(c)] La chaîne est-elle apériodique ?
    \item[(d)] Déterminer une distribution stationnaire.
\end{enumerate}}

\subsection{Algorithmique}
\exercise{1}{Suites récurrentes linéaires}{On considère l'algorithme suivant :

\begin{quote}
\texttt{def F(n) :}\\
\indent \texttt{\quad return 0 if n < 2 else 1 if n == 2 else 2 * F(n-1) + F(n-3)}
\end{quote}

\begin{enumerate}
    \item Soit $A(n)$ le nombre d'opérations arithmétiques sur des entiers effectuées lors de l'exécution de \texttt{F(n)}. Montrer que $A(n) \in \Omega(2^{n/3})$.
    \item Proposer un algorithme \texttt{Fd(n)} calculant la même valeur que \texttt{F(n)} en effectuant un nombre linéaire d'opérations arithmétiques sur des entiers, avec une complexité en espace linéaire. Quelle est sa complexité en temps ?
    \item Proposer un algorithme \texttt{Fe(n)} calculant la même valeur de manière encore plus efficace. Quelle est sa complexité en temps ?
\end{enumerate}
}

\exercise{2}{Suites récurrentes linéaires}{On considère l'algorithme suivant :

\begin{quote}
\texttt{def F(n) :}\\
\indent \texttt{\quad return 0 if n < 3 else 1 if n == 3 else 2 * F(n-1) + 4 * F(n-3) + F(n-4)}
\end{quote}

\begin{enumerate}
    \item Soit $A(n)$ le nombre d'opérations arithmétiques sur des entiers effectuées lors de l'exécution de \texttt{F(n)}. Montrer que $A(n) \in \Omega(3^{n/4})$.
    \item Proposer un algorithme \texttt{Fd(n)} calculant la même valeur que \texttt{F(n)} en effectuant un nombre linéaire d'opérations arithmétiques sur des entiers, avec la meilleure complexité en espace possible. Quelle est sa complexité en espace ? en temps ?
    \item Proposer un algorithme \texttt{Fe(n)} calculant la même valeur de manière encore plus efficace. Quelle est sa complexité en temps ?
\end{enumerate}
}
\subsection{Langage C}
Révision et entraînement écriture/lecture de fichiers

\newpage
\section{Jour 2}

\subsection{Analyse}
\exercise{1}{Produit de convolution}{Soient $f$ et $g$ deux fonctions définies sur $\mathbb{R}$, continues et $T$-périodiques ($T$ réel strictement positif). Pour $x \in \mathbb{R}$, on pose
\[ (f * g)(x) = \int_0^T f(x - t)g(t) dt. \]
\begin{enumerate}
    \item Montrer que la fonction $f * g$ est définie sur $\mathbb{R}$, continue et $T$-périodique.
    \item Montrer que $f * g = g * f$.
\end{enumerate}}

\exercise{2}{Intégrale à paramètre}{On pose, pour $x \in \mathbb{R}$,
\[ F(x) = \int_0^{+\infty} \frac{\sin(xt)}{t} e^{-t} dt. \]
\begin{enumerate}
    \item Justifier que $F$ est bien définie sur $\mathbb{R}$.
    \item Montrer que la fonction $F$ est de classe $C^1$. Pour tout nombre réel $x$, calculer $F'(x)$.
    \item En déduire une expression explicite de $F(x)$.
\end{enumerate}}

\exercise{3}{Calcul d'une intégrale double}{On pose :
\[ D = \left\{(x, y) \in \mathbb{R}^2; \ 0 \le x \le 1, \ 0 \le y \le 1, \ x^2 + y^2 \ge 1\right\}. \]
Calculer
\[ \iint_D \frac{xy}{1 + x^2 + y^2} dxdy. \]}

\exercise{4}{Application : calcul d'une intégrale impropre}{
\begin{enumerate}
    \item Calculer $\iint_{D_R} e^{-(x^2+y^2)} dxdy$, où $D_R$ est le quart de disque supérieur droit de centre $(0,0)$ et de rayon $R > 0$.
    \item Montrer que
    \[ \iint_{D_R} e^{-(x^2+y^2)} dxdy \le \iint_{[0,R]^2} e^{-(x^2+y^2)} dxdy \le \iint_{D_{\sqrt{2}R}} e^{-(x^2+y^2)} dxdy. \]
    \item En déduire la valeur de $\int_{-\infty}^{\infty} e^{-x^2} dx$.
\end{enumerate}}

\subsection{Langage C}
\exercise{1}{Révision du cours 9}{}
\exercise{2}{Copier-Couper}{
\textbf{1. Copier}
Écrire une fonction \texttt{char *copier(char *s, int pos, int nbr)}. Cette fonction doit allouer, construire, puis renvoyer l'adresse de départ d'une nouvelle chaîne contenant tous les caractères de \texttt{s} dont les positions sont comprises entre \texttt{pos} et \texttt{pos + (nbr - 1)} inclus, dans l'ordre où ils apparaissent dans \texttt{s}. Noter que ces conventions autorisent \texttt{pos} à être négatif ou au-delà de la dernière position de \texttt{s}, ou encore, \texttt{nbr} à être nul ou négatif. Parmi toutes les positions spécifiées par \texttt{pos} et \texttt{nbr}, on ignorera simplement toutes celles qui ne sont pas des positions dans \texttt{s}. Lorsque par exemple \texttt{s} contient \texttt{abcde} :

\begin{itemize}
    \item \texttt{copier(s, 2, 2)} renvoie \texttt{"cd"}, (positions 2,3)
    \item \texttt{copier(s, 2, -3)} renvoie \texttt{"abc"}, (2,1,0)
    \item \texttt{copier(s, 2, 5)} renvoie \texttt{"cde"}, (2,3,4 puis 5,6 ignorées)
    \item \texttt{copier(s, -2, 4)} renvoie \texttt{"ab"}, (-2,-1 ignorées, puis 0,1)
    \item \texttt{copier(s, 2, 0)} renvoie \texttt{""}, (est non un pointeur nul)
\end{itemize}

\textbf{2. Couper}
Écrire une fonction \texttt{char *couper(char *s, int pos, int nbr)}. Cette fonction doit allouer, construire, puis renvoyer l'adresse de départ d'une nouvelle chaîne contenant tous les caractères de \texttt{s} dans l'ordre où ils apparaissent dans \texttt{s}, \textit{sauf} ceux dont les positions sont comprises entre \texttt{pos} et \texttt{pos + (nbr-1)} inclus.
}

\subsection{Probabilités}
\exercise{1}{Marches aléatoires et loi de Rademacher}{Soit $n \ge 1$ un entier. Soient $X_1, \dots, X_n, X'_1, \dots, X'_n$ des variables aléatoires indépendantes de loi de Rademacher, c'est-à-dire,
\[ \mathbb{P}(X_i = 1) = \mathbb{P}(X_i = -1) = \mathbb{P}(X'_i = 1) = \mathbb{P}(X'_i = -1) = 1/2, \quad \forall i \in \{1, \dots, n\}. \]
On pose $S_0 = 0$, $S'_0 = 0$ et pour tout $k \in \{1, \dots, n\}$
\[ S_k = X_1 + \dots + X_k, \quad S'_n = X'_1 + \dots + X'_k. \]
On rappelle que pour tout $\ell \in \{-n, -n + 2, -n + 4, \dots, n - 2, n\}$
\[ \mathbb{P}(S_n = \ell) = 2^{-n} \binom{n}{\frac{n+\ell}{2}}. \]

\begin{enumerate}
    \item[(a)] Déterminer $\mathbb{E}[S_n]$ et $\text{Var}(S_n)$.
    \item[(b)] Déterminer la loi jointe de $(X_1 - X'_1, X_1 + X'_1)$.
    \item[(c)] Déterminer les lois de $U_n := S_n - S'_n$ et $V_n := S_n + S'_n$.
    \item[(d)] Calculer $\text{Cov}(U_n, V_n)$.
    \item[(e)] Les variables aléatoires $U_n$ et $V_n$ sont-elles indépendantes ?
\end{enumerate}}

\newpage
\section{Jour 3}
\subsection{Langage C}
\exercise{1}{Révision du cours 10}{}
\exercise{2}{Révision du cours 11}{}
\exercise{3}{Chaînes et itérations}{Dans la fonction donnée et la fonction demandée, on supposera que chaque ligne des textes considérés, se termine par un retour à la ligne, c'est-à-dire par un \texttt{\textbackslash n} - tenez en compte. Nous supposons déclaré une fonction \texttt{void afficher\_page(int k, char *s)} (vous n'avez pas à l'écrire), cette fonction affiche à l'écran le texte stocké à l'adresse \texttt{s} sur une largeur maximale \texttt{k+1} colonnes - les lignes trop longue sont coupées par autant de retours à la ligne nécessaire. Par exemple, si \texttt{k} vaut 1 et si le texte est \texttt{"abcde\textbackslash nfg\textbackslash nh\textbackslash n"} le texte affiché sera :
\begin{quote}
\texttt{ab}\\
\texttt{cd}\\
\texttt{e}\\
\texttt{fg}\\
\texttt{h}
\end{quote}
Noter que les seuls retours à la ligne ajoutés sont après \texttt{"ab"} et \texttt{"cd"} : tous les autres font parti du texte.

Écrire une fonction \texttt{char *mettre\_en\_page(int k, char *s)}. Cette fonction doit allouer (par \texttt{malloc}) et construire une chaîne dont le texte sera celui stocké à l'adresse \texttt{s} avec tous les retours à la ligne ajoutés par la fonction \texttt{void afficher\_page(int k, char *s)}. Par exemple, si $k$ vaut 1 et si le texte est \texttt{"abcde\textbackslash nfg\textbackslash nh\textbackslash n"} le texte produit sera \texttt{"ab\textbackslash ncd\textbackslash ne\textbackslash nfg\textbackslash nh\textbackslash n"}. Le texte stocké à l'adresse \texttt{s} ne devra être lu que deux fois en tout : une fois pour déterminer la taille d'allocation, une fois pour la construction du nouveau texte.}

\subsection{Analyse}

\exercise{1}{Convergence de séries de fonctions}{Étudier les convergences simple, uniforme et normale de la série de fonctions $\sum f_n$ dans chacun des cas suivants :
\begin{enumerate}
    \item[a)] $f_n(x) = nx^2 e^{-x\sqrt{n}}$ sur $\mathbb{R}^+$
    \item[b)] $f_n(x) = \frac{1}{n + n^3 x^2}$ sur $\mathbb{R}^*_+$
    \item[c)] $f_n(x) = (-1)^n \frac{x}{(1 + x^2)^n}$ sur $\mathbb{R}$
    \item[d)] $f_n(x) = \frac{\sin(n^3 x^2)}{1 + 3nx^2 + 2n^2}$ pour $x \in \mathbb{R}$
    \item[e)] $f_n(x) = nx \exp(-nx)$ pour $x \in [0, +\infty[$
\end{enumerate}}

\exercise{2}{Séries et intégrales}{Soit $a \in ]-1, 1[$. Pour $n \in \mathbb{N}$ et $t \in [0, \pi/2]$, on pose :
\[ u_n(t) = a^n \cos^n(t). \]
\begin{enumerate}
    \item Montrer que la série de fonctions $\sum u_n$ converge uniformément sur $[0, \pi/2]$. Déterminer la somme de cette série.
    \item En déduire l'égalité suivante :
    \[ \int_0^{\pi/2} \frac{dt}{1 - a \cos t} = \sum_{n=0}^{\infty} \left( \int_0^{\pi/2} \cos^n(t) dt \right) a^n. \]
\end{enumerate}}

\exercise{3}{Étude d'une fonction définie par une série}{Pour tout entier $n > 0$ et tout réel $x$ on pose :
\[ f_n(x) = \frac{\arctan(nx)}{n^2}, \]
et on considère $f(x) = \sum_{n \ge 1} f_n(x)$.
\begin{enumerate}
    \item Déterminer le domaine de définition de $f$ et sa parité éventuelle.
    \item Étudier la convergence normale de $f$ sur $\mathbb{R}$, préciser si $f$ est continue et quelles sont ses limites en $+\infty$ et en $-\infty$.
    \item Montrer que $f$ est de classe $C^1$ sur $\mathbb{R}^*$. Montrer que $f'$ est décroissante sur $]0, +\infty[$.
    \item Montrer que $f$ n'est pas dérivable en 0.
\end{enumerate}}

\newpage
\section{Jour 4}
Sera complété ultérieurement \textit{(les nouveaux jours seront ajoutés le 13 mai 2026)}.

\newpage
\section{Jour 5}
Sera complété ultérieurement \textit{(les nouveaux jours seront ajoutés le 13 mai 2026)}.

\newpage
\section{Jour 6}
Sera complété ultérieurement \textit{(les nouveaux jours seront ajoutés le 13 mai 2026)}.

\newpage
\section{Jour 7}
Sera complété ultérieurement \textit{(les nouveaux jours seront ajoutés le 13 mai 2026)}.

\newpage
\section{Jour 8}
Sera complété ultérieurement \textit{(les nouveaux jours seront ajoutés le 13 mai 2026)}.

\newpage
\section{Jour 9}
Sera complété ultérieurement \textit{(les nouveaux jours seront ajoutés le 13 mai 2026)}.

\newpage
\section{Jour 10}
Sera complété ultérieurement \textit{(les nouveaux jours seront ajoutés le 13 mai 2026)}.

\end{document}
